\documentclass[12pt]{article}
\usepackage{amsmath, amssymb, enumitem, fancyhdr, graphicx, libertine, nth,
  pgfplots, tikz, xcolor}
\pgfplotsset{compat=1.11}
\usetikzlibrary{decorations.pathreplacing}
\usetikzlibrary{patterns}
\usepackage[margin=1in]{geometry}
\usepackage[libertine]{newtxmath}
\pagestyle{fancy}
\definecolor{solution-color}{HTML}{000080}
\chead{\emph{EC201 -- Problem Set 4}}

% Define new topic command
\newcommand{\topic}[1]{\medskip\begin{center}\large\textsc{#1}\normalsize
  \end{center}}

% Define question counter:
\newcounter{question}[section]
\newcommand{\question}[1]{\refstepcounter{question}
  \subsubsection*{Question~\thequestion~-- #1}}

% Subquestions are enumerate with letters
\newenvironment{subquestions}{
  \begin{enumerate}[label=(\roman*)]}{\end{enumerate}}

\begin{document}

\topic{Elasticities}
\question{Cobb-Douglas Elasticities}
Your utility function for goods 1 and 2 is given by:
\[
  u\left( x_1,x_2 \right)=x_1^{\frac{1}{3}}x_2^{\frac{2}{3}}
\]
\begin{subquestions}
  \item Derive the demand functions $x_1\left( p_1,p_2,m \right)$ and $x_2\left(
    p_1,p_2,m \right)$.
  \item Find the own-price elasticity of demand for good 1, $\varepsilon_1$.
  \item Is good 1 an ordinary or Giffen good?
  \item Find the the cross-price elasticity of good 1 with respect to good 2,
    $\varepsilon_{12}$.
  \item Are the substitutes, complements, or neither?
  \item Find the income elasticity of demand for good 1, $\eta_1$.
  \item Is the good normal or inferior?
\end{subquestions}
When finding the elasticities, do not just write down the answer. Show each
step in the calculation.

\question{Linear demand and Marginal Revenue}
The market demand curve for a good is given by:
\[
  D\left( p \right)=10-p
\]
\begin{subquestions}
  \item What is the own-price elasticity of demand at the following price and
    quantity pairs:
    \begin{enumerate}
      \item $p=4$ and $q=6$.
      \item $p=5$ and $q=5$.
      \item $p=6$ and $q=4$.
    \end{enumerate}
  Comment on whether demand is elastic, inelastic or unit elastic at each of
  these prices.
  \item Find the inverse demand curve, $p\left( q \right)$ for the demand
    curve above.
  \item Find the revenue function, $R\left( q \right)$ for a single firm
    selling to this market.
  \item Find the marginal revenue function $MR\left( q \right)$.
  \item What is the marginal revenue at each of the price and quantity pairs in
    (i) above?
  \item What does your finding in (v) tell us about the relationship between
    elasticity and marginal revenue?
\end{subquestions}
% \question{Cobb-Douglas and Marginal Revenue}
% What is the marginal revenue of a single firm selling a good to consumers who
% all have the same Cobb-Douglas utility function?

\newpage
\topic{Monopoly}
\question{Monopoly}
The (inverse) market demand for a good is given by $p\left( q \right)=130-q$
and there is a single producer with a cost function, $c(q) = 1600+10q+q^2$.
Find the equilibrium price, quantity, and profits for the monopolist.

\question{Monopoly vs.~Perfect Competition with Linear Demand and Costs}
The inverse demand curve for a good is given by:
\[
  p\left( q \right)=10-q
\]
The cost function for the monopolist is:
\[
  c\left(  q\right)=2q
\]
\begin{subquestions}
  \item Derive the optimal price an quantity the monopolist will choose.
  \item Confirm that the monopolist charges a mark-up over marginal cost of
    $\frac{1}{1-\frac{1}{\left| \varepsilon \right|}}$. You can
    use your answer in Q2 (i).
  \item Find the profit the monopolist makes.
  \item Draw a diagram showing the following:
    \begin{itemize}
      \item The inverse demand curve.
      \item The marginal revenue curve.
      \item The marginal cost curve.
      \item The monopolist's price and quantity.
      \item The consumer surplus.
      \item The monopolist's profits.
      \item The deadweight loss.
    \end{itemize}
    The scaling of the diagram does not need to be exact. However, the
    lines/curves should have the correct shape.
  \item What is the consumer surplus?
  \item What is the deadweight loss, if any?
\end{subquestions}
Now assume the market was instead served by a large number of \textbf{perfectly
competitive} firms.
\begin{subquestions}
  \setcounter{enumi}{5}
  \item What would the price and quantity be? What profit would each of the
    firms make?
  \item Draw a diagram showing the following:
    \begin{itemize}
      \item The inverse demand curve.
      \item The marginal cost curve.
      \item The perfectly competitive price and quantity.
      \item The consumer surplus.
      \item The producer surplus.
    \end{itemize}
  \item What is the consumer surplus?
  \item What is the producer surplus?
  \item What is the deadweight loss, if any?
  \item How do each of price, quantity, profits, producer surplus, consumer
    surplus and deadweight loss compare between monopoly and perfect competition?
\end{subquestions}

\question{Natural Monopoly}
The inverse demand curve for a good is given by:
\[
  p\left( q \right)=10-q
\]
The cost function for the monopolist is:
\[
  c\left(  q\right)=9
\]
\begin{subquestions}
  \item Derive the optimal price an quantity the monopolist will choose.
  \item Find the profit the monopolist makes.
  \item What is the consumer surplus?
\end{subquestions}
Suppose the government introduced some regulation which forced the monopolist
to charge at marginal cost.
\begin{subquestions}
  \setcounter{enumi}{3}
  \item What would the new price, quantity and profits be?
  \item What is the consumer surplus?
\end{subquestions}
Suppose the government introduced some regulation which forced the monopolist
to charge a price such that it makes zero profits overall.
\begin{subquestions}
  \setcounter{enumi}{5}
  \item What would the price, quantity and profits be?
  \item What is the consumer surplus?
\end{subquestions}
There are three possible policies the government can pursue, the outcomes of
which you have solved for above:
    \begin{itemize}
      \item No regulation.
      \item Forcing the monopolist to charge at marginal cost.
      \item Forcing the monopolist to charge a price such that it makes zero
        profits.
    \end{itemize}
\begin{subquestions}
  \setcounter{enumi}{7}
  \item Which policy maximizes consumer surplus?
  \item Which policy maximizes consumer surplus \emph{plus} monopolist profits?
\end{subquestions}

\question{Two-Period Monopoly with an Experience Good}
Consider the following model of a two-period monopolist selling an
\emph{experience good}. The inverse market demand in the first period is
\[
  p\left( q_1 \right) =10-q_1
\]
In the second period the demand is:
\[
  p\left( q_2 \right) = 10+\frac{q_1}{2}-q_2
\]
The amount the monopolist sells in the second period depends on how much the
monopolist sold in the first period. The amount the monopolist sold in the
first period shifts out the second period demand curve (raises the vertical
intercept). The monopolist's cost function is $c\left( q_1,q_2 \right)=0$.
\begin{subquestions}
  \item If the monopolist was unaware he was selling an experience good and
    simply maximized profits in each period, what would his total profits be?
  \item If the monopolist knows the way the demand curves depend on both
    quantities, write down his profit function for both periods. This should be
    a function of $q_1$ and $q_2$ only.
  \item Solve for the optimal prices and quantities by maximizing the profit
    function you wrote down above. Remember that maximizing a function of two
    variables requires setting both partial derivatives equal to 0 and solving
    the system of equations. Does the monopolist offer a low introductory
    price?
  \item Compare the profits from parts (i) and (iii). Is the monopolist better
    off in (i) or (iii)?
\end{subquestions}

% \newpage
% \topic{Monopoly Behavior}
% \question{Third-Degree Price Discrimination}
% There is one movie theater in a town. The town has two groups of people:
% students and non-students. The non-students have a demand curve:
% \[
%   D_1\left( p_1\right) = 10-p_1
% \]
% and the students have a demand curve:
% \[
%   D_2\left( p_2 \right) = 10 - 2p_2
% \]
% The movie theater has a marginal cost of \$2 per customer (cleaning up their
% spilled soda and popcorn). The fixed cost is 10.
% \begin{subquestions}
%   \item Suppose first the monopolist movie theatre cannot price discriminate.
%     What price should they charge and what quantity will it sell? What
%     profit does it make?
%   \item Suppose now the monopolist is able to price discriminate. What price
%     should it charge students and what price should it charge non-students?
%     How many tickets will they sell?
% \end{subquestions}
\end{document}
